% Nouveau EGA V (cote 25): archive cote 25, site page 66 (PDF page 67).
% typescript. Draft compared: none.
% AI-transcribed by Claude Opus 5.5 (claude-opus-5-5, adaptive thinking, effort high), Message Batches API, batch ega25b (msgbatch_01A7n5YuXubUPNK1zsepi8RG), 2026-11-13; not yet checked by eye.
\scanpage{25/66}
\mspage{\hand{\unsure{21}}}

12. \uline{Généralisation de résultats antérieurs aux sections linéaires}.

Compléments de notations. Si $P = P(\underline{E})$, $\underline{E}$ Module quasi-cohérent
quelconque, on posera aussi $\mathrm{Grass}_n(P) = \mathrm{Grass}_{n+1}(\underline{E})$ , de sorte que
$\mathrm{Grass}_n(P)$ correspond aux sous-variétés linéaires de \uline{dimension} n dans
P ; ceci vaut pour $n \geq -1$ , en convenant que dim $-1$ signifie vide
Lorsque E est localement libre, \xout{\unsure{de rang}} il y a lieu également d'intro-%
duire $\mathrm{Grass}^n(P) = \mathrm{Grass}_{n-1}(\check{P}) = \mathrm{Grass}_n(\check{\underline{E}})$ , qui correspond aux sous-%
variétés linéaires de codimension n dans P. Si $\underline{E}$ est de rang r+1 i.e.
P de dimension relative r, on a un isomorphisme canonique
$\mathrm{Grass}^n(P) = \mathrm{Grass}_{r-n}(P)$ . Par la suite, nous supposons fixé E loc
libre de rang r , et nous nous intéressons aux sous-variétés linéaires
de P de codimension donnée m , donc à $\mathrm{Gr}^m = \mathrm{Gr}^m(P) = \mathrm{Gr}_{\unsure{m+1}}(\check{\underline{E}})$ .
Sur ce préschéma, il y a donc un quotient \inserted{$\underline{G}$} loc libre de rang m canoniqu
de $\underline{E}_{\mathrm{Gr}}$ , soit F . \xout{\unsure{La relation di}} Le préschéma d'incidence naturel
sur $P\times_S \mathrm{Gr}^m$ , qui représente le sous-foncteur du foncteur produit
correspond à des couples d'une section de $P_S$, et une sous-variété
linéaire de codimension m de $P_S$, \uline{contenant} cette dernière, s'explicite
alors de la façon suivante: Soit $T=P\times_S \mathrm{Gr}^m$ (ou si on préfère, un
préschéma quelconque au dessus de ce produit), alors sur \xout{E} T on a
$\underline{E}_T$ , le quotient $\underline{O}_T(1)$ , et le sous-Module localement facteur direct
$\check{\underline{G}}_T$ . Considérons l'homomorphisme composé des hom canoniques
\[ \check{G}_T \longrightarrow \underline{O}_T(1) \]
qui par transposition correspond aussi à l'homomorphisme composé
analogue \unsure{du|da} sous-Module $\underline{O}_T(-1)$ de $\check{\underline{E}}_T$ dans le quotient $\underline{G}_T{}^{\vee}$:
\[ \underline{O}_T(-1) \longrightarrow G_T \quad , \]
et peut aussi s'envisager comme défini par une section de $\underline{G}_T(1)$
\[ \phi^m \in \Gamma(T, \underline{G}_T(1)) \]

\mspage{\hand{66}}
